\section{O Índice de Qualidade Parlamentar (IQP)}

A avaliação da qualidade da representação política é uma tarefa complexa que transcende a simples análise de votações ou da produção legislativa. Um mandato parlamentar envolve múltiplas dimensões, desde a capacidade de atuação nas estruturas de poder do Legislativo até a forma como o representante se comunica com a sociedade e reflete sua composição demográfica. Com o objetivo de oferecer uma análise mais holística e multifacetada do desempenho parlamentar, propomos o \textbf{Índice de Qualidade Parlamentar (IQP)}, um indicador composto que agrega quatro dimensões distintas e complementares da atuação dos deputados federais brasileiros.

O IQP foi concebido para fornecer uma visão integrada, estruturando-se a partir dos seguintes pilares:

\begin{itemize}
    \item \textbf{Índice de Profissionalização do Parlamentar (IPP):} Avalia a trajetória e a especialização do parlamentar dentro do Congresso Nacional. Esta dimensão mensura o acúmulo de experiência e a influência institucional, considerando variáveis relacionadas à carreira legislativa e à ocupação de posições estratégicas, como a participação em mesas diretoras, lideranças e comissões.

    \item \textbf{Índice de Representação Descritiva (IRD):} Mede o grau de alinhamento entre o perfil demográfico do parlamento e o da população brasileira. Focado em características como gênero e raça, o IRD quantifica quão bem a Câmara dos Deputados espelha a diversidade da sociedade que representa, um elemento fundamental para a legitimidade do sistema político.

    \item \textbf{Índice de Representação Substantiva (IRS):} Analisa a coerência das pautas defendidas pelo parlamentar em suas comunicações públicas. Por meio da análise semântica de postagens em redes sociais, este índice mensura a consistência temática do discurso, indicando o quão focado e estável é o deputado na defesa de determinadas causas e grupos sociais.

    \item \textbf{Índice de Falar Sincero (IFS):} (em desenvolvimento) Afere o compromisso do parlamentar com a veracidade no debate público. O IFS mede a frequência com que os deputados disseminam conteúdos associados a narrativas de desinformação, penalizando aqueles que mais contribuem para a poluição do ambiente informacional.
\end{itemize}

Ao combinar essas quatro dimensões — profissionalização, representatividade demográfica, coerência discursiva e integridade informacional — o IQP oferece uma ferramenta robusta para uma análise mais rica e detalhada da qualidade parlamentar. As seções a seguir apresentam em detalhes a metodologia de cálculo de cada um dos índices que compõem o IQP.

Este documento tem como objetivo apresentar de forma detalhada a metodologia por trás do IQP. A exposição está organizada da seguinte forma: inicialmente, são descritos os quatro pilares que compõem o índice, explicitando a relevância de cada um para a avaliação da qualidade parlamentar. Em seguida, as seções subsequentes são dedicadas a aprofundar a construção de cada indicador individualmente: o Índice de Profissionalização do Parlamentar (IPP), o Índice de Representação Descritiva (IRD), o Índice de Representação Substantiva (IRS) e, por fim, o Índice de Falar Sincero (IFS). Cada seção abordará as fontes de dados, as variáveis utilizadas e os procedimentos de cálculo específicos.
